\chapter{Computer generation of numerical and algebraic expressions}\label{chap:14}\index{computer algebra}\index{computer generation of tables}

The algebraic tables of Secs.~\ref{sec:8.12}, \ref{sec:9.11} and \ref{sec:10.11},
and the numerical tables of Secs.~\ref{sec:8.13} and \ref{sec:10.12}, together
with Tables 9.9--9.11, may be reproduced entry by entry by computer algebra. This
chapter records the method used, so that any individual entry can be
regenerated and checked rather than taken on trust. The programs are listed in
Sec.~\ref{sec:14.5}.

\section{Scope}\label{sec:14.1}\index{computer algebra!scope}
Each algebraic table fixes one small angular momentum and leaves the remaining
arguments symbolic: $b$ in the Clebsch-Gordan tables, $d$ in the $6j$ tables,
and the triad $(\alpha \beta \gamma)$ in the $9j$ tables. The corresponding sum
is then finite with a number of terms bounded by that fixed argument, and is
therefore a finite symbolic computation. Nothing in the method is restricted to
the ranges printed here; the same programs generate entries beyond them, at
greater cost.

\section{Algebraic generation}\label{sec:14.2}\index{computer algebra!algebraic tables}

\subsection{The Clebsch-Gordan coefficients}\label{sec:14.2.1}\index{Clebsch-Gordan coefficient!computer generation}
For fixed $b$ and $\beta$, and with $c=a+k$, the summation index of the Van der
Waerden-Racah sum of Sec.~\ref{sec:8.2.1} runs over at most $2b+1$ values, so
the coefficient is a short symbolic sum. The difficulty is that its terms
contain factorials of $a\pm\alpha$, which are not polynomial in $a$. It is
removed by the change of variable
\begin{equation}
p=a+\alpha, \qquad q=a-\alpha ,\label{eq:14:2:1}
\end{equation}
both of which are non-negative \emph{integers} even when $a$ and $\alpha$ are
half-integral. Every factorial argument is then an explicit integer offset from
$p$, $q$ or $p+q=2a$, and each such factorial may be replaced by a rational
function through
\begin{equation}
\frac{(x+m)!}{x!}=
\begin{cases}
(x+1)(x+2)\cdots(x+m), & m \geq 0,\\[2pt]
\bigl[x(x-1)\cdots(x+m+1)\bigr]^{-1}, & m<0 .
\end{cases}\label{eq:14:2:2}
\end{equation}
The factor $p!\,q!$ produced by the sum cancels exactly against the one under
the square root, and the coefficient takes the form
\begin{equation}
\clebsch{a}{\alpha}{b}{\beta}{c}{\gamma}=T(p,q)\,\bigl[W(p,q)\bigr]^{1/2},
\label{eq:14:2:3}
\end{equation}
with $T$ and $W$ explicit rational functions. No general-purpose simplification
of factorials or gamma functions is required at any stage.

Eq.~\eqref{eq:14:2:3} does not fix the presentation, since any factor may be
moved under the radical as its square. The tables of Sec.~\ref{sec:8.12} absorb
every factor that is of one sign throughout the physical domain
$|\gamma| \leq c$, such as $(c+\gamma)$ or $(c-\gamma+1)$, and leave outside
only those that change sign, such as $\gamma$ itself. Applying that rule
reproduces the printed entries.

\subsection{The $6j$ symbols}\label{sec:14.2.2}\index{6j symbol@$6j$ symbol!computer generation}
The analogous integer variables for the $6j$ symbols are the triangle deficits
\begin{equation}
u=s-2a, \qquad v=s-2b, \qquad w=s-2c, \qquad s=a+b+c=u+v+w .\label{eq:14:2:4}
\end{equation}
Writing the summation index of Eq.~\eqref{eq:9:2:1} as $N=s+t$, and taking
$e=c+n$ and $f=b+m$ as in the tables, the seven factorial arguments of that
equation become
\begin{equation}
\begin{gathered}
N-a-b-c=t, \qquad N-a-e-f=t-n-m, \\
N-b-d-f=v+t-d-m, \qquad N-c-d-e=w+t-d-n, \\
a+b+d+e-N=d+n-t, \qquad a+c+d+f-N=d+m-t, \\
b+c+e+f-N=u+n+m-t ,
\end{gathered}\label{eq:14:2:5}
\end{equation}
each an explicit integer offset from $t$, $u$, $v$ or $w$. The index $t$ runs
over the explicit finite range
\begin{equation}
\max(0,\,n+m) \leq t \leq \min(d+n,\,d+m),\label{eq:14:2:6}
\end{equation}
so the sum has at most $2d+1$ terms. Here the cancellation is complete:
$\Delta(abc)$ contributes $[u!\,v!\,w!/(s+1)!]^{1/2}$, the sum contributes
$(s+1)!/(u!\,v!\,w!)$, and $\Delta(aef)$ contributes the inverse, leaving
\begin{equation}
\sixj{a}{b}{c}{d}{e}{f}=(-1)^{s}\,T(u,v,w)\,\bigl[W(u,v,w)\bigr]^{1/2}
\label{eq:14:2:7}
\end{equation}
with no symbolic factorial remaining. The phase arises because
$(-1)^{N}=(-1)^{s}(-1)^{t}$, which is why every entry of Sec.~\ref{sec:9.11}
carries $(-1)^{s}$ or $(-1)^{s+1}$.

The factors of $T$ and $W$ are linear in $a$, $b$, $c$ and are printed in the
notation of the tables by reading off their coefficients: $(1,1,1)$ gives
$s+k$, $(-1,1,1)$ gives $s-2a+k$, $(0,2,0)$ gives $2b+k$, and so on. The factor
left outside the radical is expressed as a polynomial in $X$ by division, $X$
being quadratic in $a$.

\subsection{The $9j$ symbols}\label{sec:14.2.3}\index{9j symbol@$9j$ symbol!computer generation}
The $9j$ symbols of Sec.~\ref{sec:10.11} are not given by a single sum of the
above type. They are obtained from the sum over three $6j$ symbols,
Eq.~\eqref{eq:10:2:20}. For the tabulated symbols the triad $(\beta a x)$
restricts the intermediate momentum to
\begin{equation}
x=a+\xi, \qquad -\beta \leq \xi \leq \beta ,\label{eq:14:2:8}
\end{equation}
an explicit sum of at most $2\beta+1$ terms. Each of the three $6j$ symbols is
carried by a permutation of its columns together with an interchange of the
upper and lower arguments in two of them into the form
$\bigl\{{A\ B\ C \atop d\ C{+}n\ B{+}m}\bigr\}$ treated in
Sec.~\ref{sec:14.2.2}, so the same machinery evaluates all three; only the
deficits \eqref{eq:14:2:4} need re-expressing in terms of those of $(abc)$.

The phases collapse. Summing the three factors $(-1)^{A+B+C}$ with the
$(-1)^{2x}$ of Eq.~\eqref{eq:10:2:20} gives the exponent
$4a+4\xi+2(a+b+c)+\beta+\mu$, and since $2a$, $2\xi$ and $a+b+c$ are integers
every term is even except the last two:
\begin{equation}
\ninej{a+\lambda}{b+\mu}{c+\nu}{a}{b}{c}{\alpha}{\beta}{\gamma}
=(-1)^{\beta+\mu}\sum_{\xi}(2a+2\xi+1)\,T_1T_2T_3\,[W_1W_2W_3]^{1/2},
\label{eq:14:2:9}
\end{equation}
independent of $\xi$ and of $a$, $b$, $c$. The radicands of the individual
terms differ by perfect squares of rational functions, so a common radical may
be extracted and the whole reduces to a rational multiple of a single square
root. That the ratios are perfect squares is tested, not assumed.

\section{Numerical evaluation}\label{sec:14.3}\index{computer algebra!numerical tables}
Numerical values are obtained in exact rational arithmetic. Since every
coefficient of Secs.~\ref{sec:8.13} and \ref{sec:10.12}, and of
Tables 9.9--9.11, is a signed square root of a rational number, the value $V$ is squared, $V^{2}$ is
confirmed to be rational, and the entry is presented as
\begin{equation}
V=\pm\left(\frac{P}{Q}\right)^{1/2},\label{eq:14:3:1}
\end{equation}
the decimal being derived from the fraction rather than the reverse. The
$12j$ symbols are evaluated from their defining single sums over four $6j$
symbols, Eqs.~\eqref{eq:10:13:6} and \eqref{eq:10:13:26}.

\section{Verification}\label{sec:14.4}\index{computer algebra!verification}
Each generator is checked against an independent implementation of the same
coefficient at a large number of numerical points, at two levels: the sum
itself, and the regrouping into the printed form, so that an error in a phase
or in the choice of factors placed under the radical cannot pass unnoticed.
Points are chosen asymmetrically as well as symmetrically, since a factor such
as $X$ or $Z$ changes sign only when one momentum dominates the other two, and
a symmetric sample would not detect a factor wrongly treated as of one sign.
All entries of the tables cited in Sec.~\ref{sec:14.1} agree. The relations of
Sec.~\ref{sec:10.13} are checked against each other in the same way.

\section{The programs}\label{sec:14.5}\index{computer algebra!programs}
The generators and their checks are written in Python using
\textsf{SymPy}. Each generator produces a single entry, given the fixed
argument and the row and column indices of the table:
\begin{lstlisting}[language=bash]
gen_8_12_cg_tables.py  --b 1 --beta=0 --k=0
gen_9_11_6j_tables.py  --d 1 --m=0 --n=0
gen_10_11_9j_tables.py --alpha 3/2 --beta 3/2 --gamma 0 --lam=1/2 --mu=1/2 --nu=0
\end{lstlisting}
Omitting the row and column indices produces a whole table. The companion
programs \texttt{check\_8\_12.py}, \texttt{check\_9\_11.py},
\texttt{check\_10\_11.py} and \texttt{check\_10\_13.py} perform the
verification of Sec.~\ref{sec:14.4}. The same functions are available
interactively in the two notebooks described in Sec.~\ref{sec:14.6}.

\section{Use of the notebooks}\label{sec:14.6}\index{computer algebra!notebooks}
Two Jupyter notebooks make the generators available without recourse to the
command line. Both require only \textsf{SymPy}; if \textsf{ipywidgets} is
also present the arguments are set from menus, and if it is not, each quantity
is instead obtained by editing the arguments at the head of a cell and running
it. Every result is displayed both as a typeset formula and as the
corresponding \textsf{SymPy} expression, the latter being the form to copy if
the entry is to be used in further computation.

\subsection{Algebraic expressions}\label{sec:14.6.1}
The notebook \texttt{vmk\_algebraic\_tables.ipynb} has one part for each of
the three families of Sec.~\ref{sec:14.2}. In each, one selects the fixed
argument and the row and column of the table:
\begin{lstlisting}[language=Python]
clebsch(b=1, beta=0, k=0)             # <a alpha, b beta | c gamma>, c = a+k
sixj(d=1, m=0, n=0)                   # {a b c; d e f}, e = c+n, f = b+m
ninej(1, 1, 0, lam=1, mu=0, nu=0)     # alpha, beta, gamma, then lambda, mu, nu
\end{lstlisting}
A further cell in each part prints an entire table at once. Setting
\texttt{show\_source=True} adds the \LaTeX{} source of the entry. The
Clebsch-Gordan entries are recomputed as soon as an argument is changed; the
$6j$ and $9j$ entries are not, since a single entry of the larger tables takes
appreciably longer to construct, and are instead computed on request. A closing
part runs the verification programs of Sec.~\ref{sec:14.4} over a reduced
range of arguments.

\subsection{Numerical values}\label{sec:14.6.2}
The notebook \texttt{vmk\_numeric\_tables.ipynb} evaluates the
Clebsch-Gordan coefficients and the $6j$, $9j$ and $12j$ symbols at given
numerical arguments, in the form \eqref{eq:14:3:1}:
\begin{lstlisting}[language=Python]
clebsch(1, 0, 1, 0, 2)                # gamma = alpha + beta is supplied
sixj(1, 1, 1, 1, 1, 1)
ninej(0.5, 0.5, 1, 0.5, 0.5, 1, 1, 1, 2)
twelvej1(1, 1, 1, 1, 1, 1, 1, 1, 1, 1, 1, 1)    # eq (10.13.6)
twelvej2(1, 1, 1, 1, 1, 1, 1, 1, 1, 1, 1, 1)    # eq (10.13.26)
\end{lstlisting}
A coefficient that vanishes is reported together with the condition that fails,
whether a triangular condition on a particular triad or a projection outside
its permitted range, so that a zero is never returned without explanation.
Further cells reproduce the numerical tables themselves under the conditions
stated there, and a final part evaluates the algebraic expressions of
Sec.~\ref{sec:14.2} at numerical arguments and compares them with the values
obtained directly, the two being independent routes to the same quantity.
